\documentclass[12pt,a4paper]{article}

% Page layout
\usepackage[
    margin=2.5cm
]{geometry}

% Font and encoding
\usepackage[T1]{fontenc}
\usepackage[utf8]{inputenc}
\usepackage{lmodern}

% Line spacing
\usepackage{setspace}
\doublespacing

% Line numbers
\usepackage[left]{lineno}

% Mathematics
\usepackage{amsmath}
\usepackage{amssymb}

% Figures
\usepackage{graphicx}

% Tables
\usepackage{booktabs}

% Units
\usepackage{siunitx}
\DeclareSIUnit{\pixel}{px}
\DeclareSIUnit{\byte}{B}

% References
\usepackage[
    round,
    authoryear
]{natbib}

% Links
\usepackage[
    colorlinks=true,
    citecolor=blue,
    linkcolor=blue,
    urlcolor=blue
]{hyperref}


\title{
    Mapping of Post-Fire Shrub Vegetation Using Deep Learning and UAV RGB-LiDAR Data
}

\author{
    Cristian Iranzo,
    Raquel González,
    Abdulhakim M. Abdi,
    Raquel Montorio,
    Fernando Pérez-Cabello
}

\date{}

\begin{document}

\maketitle

% Start line numbering
\linenumbers

\begin{abstract}
\end{abstract}

\noindent
\textbf{Keywords:}
keyword 1; keyword 2; keyword 3; keyword 4; keyword 5

\section{Introduction}

Accurate characterization of forest cover is essential for sustainable ecosystem management, land-use planning, ecological research, and the design of restoration and risk-prevention measures \citep{chazdon2008beyond, townshend2012global}. In Mediterranean ecosystems, shrublands are also particularly important because they protect soils from erosion, support biodiversity, contribute to vegetation succession, and influence fuel accumulation and wildfire behavior \citep{keeley2005determinants}.

Mediterranean shrublands consist mainly of low, densely branched woody vegetation adapted to long, dry summers, mild winters, and recurrent fire. Many shrub species recover rapidly after fire through vegetative resprouting or fire-stimulated germination, making shrubland development an important indicator of post-fire regeneration \citep{keeley1986resilience,pausas2008wildfires}. In fire-affected areas, accurately mapping shrub vegetation is therefore necessary to assess ecological recovery, estimate biomass accumulation, evaluate future fire risk, and plan restoration and forest-management actions. However, shrubland identification remains challenging because of its high spatial and spectral heterogeneity and its spectral similarity to other land-cover types \citep{liu2026underestimated,west2025landsat}.

Remote sensing has substantially improved the detection and characterization of vegetation by enabling observations over large areas at high spatial, spectral, and temporal resolutions. Traditional land-cover mapping has relied on visual interpretation, spectral classification, and manual digitization, as in regional programs such as CORINE \citep{grekousis2015overview}. Object-Based Image Analysis has also been applied by classifying imagery into homogeneous regions and assigning classes according to spectral, textural, and spatial properties \citep{blaschke2010object,soubry2022mapping}. Although these approaches can produce reliable maps, their performance depends heavily on operator expertise, parameter selection, and analysis scale, and they may have difficulty representing structurally complex and heterogeneous vegetation \citep{mishra2014mapping}.

Recent advances in artificial intelligence (AI), particularly deep learning (DL), have created new opportunities for overcoming these limitations. Convolutional Neural Networks (CNNs) can automatically extract hierarchical and multiscale features from imagery and have improved the classification and segmentation of vegetation and land-cover categories in complex environments \citep{abdi2020land,kattenborn2021review}. Architectures such as U-Net, Mask R-CNN, and DeepLab have been successfully used to delineate forest cover, shrubland, and individual tree structures from high- and very-high-resolution imagery \citep{he2023improvements,guirado2017deeplearning}.

Despite these advances, the automatic delineation of heterogeneous shrub vegetation in post-fire environments remains insufficiently explored. Addressing this gap could improve the efficiency and reliability of post-fire vegetation monitoring while reducing the time requirements and subjectivity associated with manual interpretation. However, the application of CNNs relies on substantial amounts of labeled training data, which are difficult to obtain, particularly in highly diverse ecosystems and in scenarios involving polymorphic shrubs, such as those found in post-fire environments. These conditions increase the difficulty of plant delineation in highly diverse ecosystems \citep{zhang2020identifying}, particularly for polymorphic vegetation where crown overlap and splitting may occur \citep{dong2019progressive}.

To alleviate this limitation, transfer-learning strategies can increase the applicability of CNNs by reusing features learned from previously analyzed datasets \citep{tan2018survey}. This approach reduces both computational requirements and the amount of labeled training data required, which is particularly valuable in environmental applications where reference data are often limited.

Previous studies have also shown that combining spectral, geometric, photogrammetric, and LiDAR-derived information can improve shrub and tree classification. For example, \citet{carbonellrivera2022classification} reported an overall accuracy of approximately \qty{96}{\percent} when classifying Mediterranean shrub and tree species using UAV-derived point clouds and supervised classifiers. Other studies have obtained promising results through multispectral classification, object-based analysis, and rule-based methods incorporating vegetation-height information \citep{uyeda2020vegetation,zahidi2015object}.

The aim of this study is to assess the effectiveness of very-high-spatial-resolution RGB imagery for the automated mapping of shrub cover in a fire-affected forest area. The JuniperMapper dataset \citep{khaldi2023junipermapper} is used to fine-tune the models and assess their performance in Mediterranean shrub mapping. The study comprises one Mask R-CNN implementation and one promptable transformer-based segmentation model. The Detectron2 implementation of Mask R-CNN reproduced the configuration used in the original JuniperMapper study \citep{khaldi2024individual}, providing a literature-based reference for comparison. The SAM 3 promptable model \citep{carion2026sam3} is also investigated. The classifications generated by the RGB-based models are combined with UAV LiDAR-derived point-cloud data to determine whether the integration of structural information improves shrub-mapping accuracy.

\section{Methodology}

The workflow of this research is explained in Figure~\ref{fig:methods}. First, the data were collected in the field, i.e., the drone flight and the shrub data within the locations used to validate the models. Then the data were processed, and the shrub masks were digitized. The models were trained and internally evaluated with the JuniperMapper dataset. Finally, the best configuration of each model was deployed in the study area and the metrics were computed using ground reference data. After the creation of the shrub maps, the refinement operation was applied to try to improve the masks and the accuracy metrics were recomputed. The code to reproduce the workflow is located inside the following GitHub repository: \href{https://github.com/CristianICS/shrub-segmentation}{https://github.com/CristianICS/shrub-segmentation}.

\begin{figure}[ht]
    \centering
    \includegraphics[width=1\textwidth]{figures/fig_methods.png}
    \caption{
        Phases of the methodology followed in this work. CHM stands for \textit{Canopy Height Model}.
    }
    \label{fig:methods}
\end{figure}

\subsection{Study area}

The study area is located within the Sierra de Algairén (Figure~\ref{fig:location}), which reaches elevations of approximately \qty{1000}{\metre} above sea level and is characterized by a Mediterranean climate with continental influences. The climate is characterized by cold winters and hot, dry summers, with marked temperature fluctuations throughout the year. The mean annual temperature is approximately \qty{12}{\celsius}, with winter minima that can fall to \qty{-2}{\celsius} and summer maxima frequently exceeding \qty{30}{\celsius}. Mean annual precipitation ranges between \qty{500}{\milli\metre} and \qty{600}{\milli\metre} and is concentrated mainly in spring and autumn. This precipitation regime, together with the dryness of the summer months, promotes soil desiccation and increases the risk of forest fires.

\begin{figure}[ht]
    \centering
    \includegraphics[width=1\textwidth]{figures/fig_location.png}
    \caption{
        Location map of the study area, containing the $3 \times 3$ plots used to collect the field data.
    }
    \label{fig:location}
\end{figure}

The area selected for model validation is located in the municipality of Encinacorba, Zaragoza (Aragón), within the area affected by the wildfire that occurred on 29 July 2003. A total area of \qty{267.4}{\hectare} was burned. More specifically, the study area comprises \qty{13.3}{\hectare} with a south-facing aspect that was completely burned during the wildfire (Figure~\ref{fig:location}). Its steep mean slope gradient of \qty{21.5}{\degree} has hindered post-fire vegetation recovery.

At present, the study area is dominated by low- to medium-stature shrubland. The most abundant species include \textit{Thymus vulgaris} L., \textit{Sedum sediforme} Jacq., \textit{Lavandula stoechas} L., and \textit{Cistus laurifolius} L. Scattered woody vegetation is also present, including \textit{Quercus ilex} L., \textit{Pinus pinaster} Aiton, and \textit{Crataegus monogyna} Jacq. A high proportion of bare soil and necromass remains.

\subsection{Data acquisition}

The study area was flown with a DJI Matrice 300 drone, equipped with a \qty{20}{\mega\pixel} RGB camera with a focal length of \qty{8.8}{\milli\metre}/\qty{24}{\milli\metre} integrated into the Zenmuse L1 sensor. The RTK positioning system, connected to the ARAGEA regional active geodetic network, was used to obtain the coordinates of each image with centimetre-level accuracy.

The flight mission took half an hour and started at 11:30~a.m. on 17 June 2024. The area was flown over at a constant altitude of \qty{100}{\metre} above the terrain, maintaining a constant height relative to the ground throughout the entire flight, and at a speed of \qty{3}{\metre\per\second}. Both the lateral and forward overlap were set to \qty{80}{\percent}. From the RGB images captured with the camera, an orthophoto was generated using WebODM software (v.3.2.6; \citealp{webodm}), following the default parameters for orthomosaic generation. Coordinate information and other image metadata were extracted from the EXIF data of each image. An orthomosaic with a spatial resolution of \qty{5}{\centi\metre} was obtained, along with the point cloud collected with the Zenmuse L1 LiDAR payload.

\subsection{Shrub instance-segmentation models}

\subsubsection{Dataset preparation}

The deep-learning models were developed using the JuniperMapper dataset described by \citet{khaldi2023junipermapper}. The dataset was originally designed for the individual delineation of medium- and large-sized \textit{Juniperus} shrubs and includes reference annotations derived from both expert photo-interpretation (PI) and field observations.

The PI dataset comprises 712 image tiles of 448 $\times$ 448 pixels, divided into predefined training (570), validation (67) and test (75) subsets. These datasets were employed during the two steps training of the models. First, the model weights were fine-tuned on the training subset. Neither model underwent a hyperparameter search; each used a fixed configuration taken from previous work. Second, a checkpoint selection phase was conducted on the validation subset using the highest mask AP$_{50}$.

Mask AP was computed with the standard dataset-level COCO evaluation. Mask AP$_{50}$ uses an IoU threshold of 0.50, and mask AP$_{50:95}$, which averages AP over IoU thresholds from 0.50 to 0.95 in increments of 0.05, was used as a secondary criterion to resolve ties. 

Validation mask AP was used only to select checkpoints within each architecture, not to compare architectures. Architectures were compared using the post-training evaluation with the JuniperMapper test set and the field-work (FW) dataset. The FW dataset contains 1,771 field-derived shrub masks across 124 images. It was used to select the confidence threshold of each model (Section~\ref{sec:post_training_ev}) and to assess how the fixed models transferred from photo-interpreted reference data to field-derived shrub delineations.

\subsubsection{Mask R-CNN}

The first model was a Mask R-CNN implemented in Detectron2 and configured following \citet{khaldi2024individual}. The architecture used a ResNet-101 backbone with a C4 feature-extraction head and was initialized from weights pretrained on the COCO dataset.

The model was subsequently fine-tuned end-to-end using the configuration derived from \citet{khaldi2024individual}, including a base learning rate of 0.0025, 4,000 training iterations, and an ROI-head batch size of 256 proposals per image (Table~\ref{tab:maskrcnn_hparams}). Hyperparameters not reported by \citet{khaldi2024individual} were kept at their Detectron2 default values. The model was evaluated on the PI validation subset every two epochs during training.

\begin{table}[ht]
    \centering
    \caption{
        Hyperparameters of the Mask R-CNN implementation. Values marked with $^{\dagger}$ are inherited from the Detectron2 defaults for the \texttt{mask\_rcnn\_R\_101\_C4\_3x} configuration.
    }
    \label{tab:maskrcnn_hparams}
    \begin{tabular}{ll}
        \toprule
        Hyperparameter & Value \\
        \midrule
        \multicolumn{2}{l}{\textit{Architecture}} \\
        Backbone & ResNet-101 (C4) \\
        Initial weights & COCO-pretrained (\texttt{R\_101\_C4\_3x}) \\
        Number of classes & 1 (shrub) \\
        ROI-head batch size per image & 256 \\
        \midrule
        \multicolumn{2}{l}{\textit{Optimization}} \\
        Optimizer & SGD$^{\dagger}$ \\
        Momentum & 0.9$^{\dagger}$ \\
        Weight decay & \num{1e-4}$^{\dagger}$ \\
        Base learning rate & 0.0025 \\
        Learning-rate schedule & Cosine decay to 0 \\
        Warm-up & Linear, 2,000 iterations (factor 0.1) \\
        Batch size & 2 images \\
        Training iterations & 4,000 \\
        \midrule
        \multicolumn{2}{l}{\textit{Input}} \\
        Training resize (shorter / longer edge) & 800 / 1,333 pixels \\
        Test resize (shorter / longer edge) & 800 / 1,333 pixels \\
        Data augmentation & None \\
        \midrule
        \multicolumn{2}{l}{\textit{Evaluation and checkpointing}} \\
        Evaluation and checkpoint period & 2 epochs \\
        % the script computes 2 x ceil(570 / 2) = 570 iterations, assuming 570 training images. 
        % That means the 4,000-iteration run is about 14 epochs
        Checkpoint selection criterion & Validation mask AP$_{50}$ \\
        Maximum detections per image & 256 \\
        Random seed & 42 \\
        \bottomrule
    \end{tabular}
\end{table}

\subsubsection{SAM 3}

The second model was SAM 3 \citep{carion2026sam3}, which combines prompt-conditioned detection and segmentation with a shared Perception Encoder for visual and language representations.

Two SAM 3 regimes were evaluated. First, the pretrained model was applied in zero-shot mode using the text prompt ``shrub'', without domain-specific training. This experiment provided a baseline for assessing the transferability of the pretrained representation to overhead shrub imagery.

Second, SAM 3 was adapted to the JuniperMapper domain. The vision and language backbones were kept frozen, and only the detection and segmentation components were fine-tuned on the PI training subset, using ``shrub'' as the text prompt. Training ran for 30 epochs with a batch size of 2. The trainable parameters were optimized with AdamW, using a base learning rate of \num{8e-5} with an inverse-square-root schedule, 20 warm-up steps and a final cooldown. Weight decay was 0.1, excluding biases and normalization layers, and gradients were clipped at a norm of 0.1. The objective function combined Hungarian-matched L1 and generalized-IoU bounding-box losses, a focal classification/presence loss, focal and Dice mask losses, and an auxiliary one-to-many matching term.

Data processing followed the original SAM 3 fine-tuning configuration. Training images were resized to a random size between 480 and 1008 pixels and padded to the model's 1008-pixel input resolution, and prompt boxes were jittered with Gaussian noise ($\sigma$ = 0.1, maximum displacement 20 pixels). Validation images were resized directly to 1008 pixels. The model was evaluated on the PI validation subset after each epoch.

\subsubsection{Computational environment}

Mask R-CNN and SAM 3 were trained and evaluated on a workstation equipped with an Intel Core i7 processor, \qty{64}{\giga\byte} RAM, and an NVIDIA GeForce RTX 5070 Ti GPU with \qty{16}{\giga\byte} VRAM, under Ubuntu running through WSL2 on Windows 11 with CUDA acceleration.

\subsection{Post-training evaluation}\label{sec:post_training_ev}

After model development, the selected checkpoint from each architecture was evaluated on the PI test and FW datasets. Following \citet{khaldi2024individual}, prediction performance was examined across a range of confidence-score thresholds to characterize the precision-recall trade-off and determine an appropriate operating point for deployment.

Performance was characterized using instance-level mask IoU and the directional S-IoU formulation proposed by \citet{khaldi2024individual}. For conventional IoU, predicted and reference instances were evaluated using their best-matching counterpart, and a detection was considered a true positive when the mask overlap exceeded the specified IoU threshold. IoU thresholds of 0.50 and 0.75 were considered, and the same two values were applied to S-IoU. Hereafter, the matching criterion is written as IoU~=~0.50, IoU~=~0.75, S-IoU~=~0.50 or S-IoU~=~0.75, whereas precision, recall, F1, AP and confidence thresholds are reported as percentages.

The S-IoU was configured by \citet{khaldi2024individual} to account for fragmentation and merging of shrub crowns. The metric evaluates the proportion of shrub areas being detected by the model, dividing the intersection by the ground truth area. Rather than considering only the best-matching instance, S-IoU compares each prediction or reference object with the union of the overlapping instances in the opposite set. Prediction- and reference-directed evaluations were therefore performed separately to determine true positives (TPs), false positives (FPs), and false negatives (FNs).

The Equation~\ref{eq:siou_labls} evaluates the labels and was used to compute the number of TPs and FPs. Otherwise, the Equation~\ref{eq:siou_pred} works for the predictions and identifies the number of FNs. $S_{p_{\mathrm{match}}}$ represents the set of all matching predictions and $S_{l_{\mathrm{match}}}$ all matching labels or reference masks. $l_k$ and $p_k$ are the label and prediction of id $k$, respectively.

\begin{equation}
\text{S-IoU}(p, S_{l_{\mathrm{match}}}) = \frac{
    \mathrm{Area}(p \cap \bigcup_{k=1}^{|S_{l_{\mathrm{match}}}|} l_k)
}{
    \operatorname{Area}(\bigcup_{k=1}^{|S_{l_{\mathrm{match}}}|} l_k)
}
\label{eq:siou_labls}
\end{equation}

\begin{equation}
\text{S-IoU}(l, S_{p_{\mathrm{match}}}) = \frac{
    \mathrm{Area}(l \cap \bigcup_{k=1}^{|S_{p_{\mathrm{match}}}|} p_k)
}{
    \operatorname{Area}(l)
}
\label{eq:siou_pred}
\end{equation}

S-IoU was set to 0 for objects without any overlapping counterpart. Given a matching threshold $\tau$ (0.50 or 0.75), a prediction was counted as a true positive (TP) if $\text{S-IoU}(p, S_{l_{\mathrm{match}}}) \geq \tau$ and as a false positive (FP) otherwise, and a reference mask was counted as a false negative (FN) if $\text{S-IoU}(l, S_{p_{\mathrm{match}}}) < \tau$. Counts were accumulated over all images, and precision, recall and F1-score were calculated as in Equation~\ref{eq:prf1}.

\begin{equation}
    \text{Precision} = \frac{\text{TP}}{\text{TP} + \text{FP}},
    \qquad
    \text{Recall} = \frac{\text{TP}}{\text{TP} + \text{FN}},
    \qquad
    \text{F1} = \frac{2 \cdot \text{Precision} \cdot \text{Recall}}{\text{Precision} + \text{Recall}} .
    \label{eq:prf1}
\end{equation}

Because TPs and FNs are determined from different directions, the sum of TPs and FNs need not equal the number of reference masks.

Precision, recall, and F1-score were calculated across prediction-confidence thresholds for both the PI test and FW datasets. Consistent with \citet{khaldi2024individual}, the PI test dataset was used to characterize model performance within the photo-interpretation domain, whereas the independently field-derived FW dataset was used to characterize performance under conditions more representative of real-world deployment.

For each model, F1-score response curves across confidence thresholds were calculated separately for the PI test and FW datasets. The operational confidence threshold for subsequent deployment was selected using the FW dataset, as the field-based reference data provide the closest available approximation to the conditions encountered during large-area application. True positives were defined using S-IoU~=~0.50. The operational point was defined as the highest confidence threshold for which the F1-score was within 1 percentage point of the maximum F1-score obtained on the FW dataset. This criterion favors more confident predictions and reduces false-positive detections while maintaining near-optimal F1 performance. The selected threshold and its corresponding precision, recall, and F1-score were reported together with the full confidence-threshold response curves.

\subsection{Independent study-area validation}

Field validation was conducted in the Encinacorba study area using 26 plots of \qty{9}{\square\metre} distributed along two longitudinal transects. Within each plot, vegetation composition and percentage cover were recorded. Shrub and tree objects were subsequently delineated from the field-reference information and corresponding orthomosaic imagery, producing 58 shrub polygons and 21 tree polygons.

The 58 shrubs contain 46 \textit{L. stoechas}, 8 \textit{C. laurifolius}, 4 \textit{Erica vagans} L.; and the 21 trees contain 13 \textit{Q. ilex}, 6 \textit{P. pinaster}, and 2 \textit{C. monogyna}.

\subsection{Orthomosaic inference and mask reconstruction}

Each selected model was deployed over the orthomosaic using overlapping tiles of 448 $\times$ 448 pixels with a 64-pixel overlap, matching the size of the JuniperMapper images used for training. The same tiling was used for both models. SAM 3 internally resizes every input to its fixed 1008-pixel resolution, keeping shrubs at the scale used for checkpoint selection. Tiles were passed to the models as unmodified 8-bit RGB, and pixels outside the orthomosaic footprint were excluded using its alpha band. Predicted instance masks were converted to georeferenced polygons. Instances smaller than 4 pixels (\qty{0.01}{\square\metre}) were discarded.

To reduce tile-boundary fragmentation, polygons from different inference windows that intersected, touched or were separated by no more than \qty{0.05}{\metre} were grouped and unioned, and small gaps between the merged parts were closed. Polygons from the same window were not merged. The maximum, mean and median confidence scores of each group were retained.

The resulting shrub maps were compared with the field reference. Because most reference polygons extended beyond the plots, no geometry was clipped. Each reference polygon was matched against all predictions, and a matched prediction was counted as a true positive wherever it lay. Since all shrubs within each plot were digitized, an unmatched prediction was counted as a false positive only when at least half of its area lay inside a plot; outside the plots, it could correspond to an undigitized shrub.

Predicted shrubs were matched to reference shrub polygons using greedy one-to-one matching with an IoU threshold of 0.10. Precision was calculated as the proportion of true positives among true and false positives, recall as the proportion of reference shrubs successfully matched, and F1 as the harmonic mean of precision and recall.

Tree commission was calculated as the proportion of reference tree polygons matched by a prediction, using the same one-to-one matching (IoU $\geq$ 0.10) applied separately to the tree polygons. It describes where part of the false positives fell and is not a separate test.

\subsection{SAM 3 prompt sensitivity analysis}

The text prompt ``shrub'' was defined as the primary SAM 3 prompt before any evaluation. It was the category name used to domain-adapt the SAM 3 decoder and to select the checkpoint on the validation split, and it was used for the main comparison between architectures.

A supplementary prompt-sensitivity analysis then evaluated nine alternative short noun phrases describing shrub vegetation: ``bush'', ``low shrub'', ``small shrub'', ``small bush'', ``woody shrub'', ``Mediterranean shrub'', ``dry shrub'', ``burned shrub'' and ``brown shrub''. To keep the analysis consistent with model development, each prompt was scored with the same data and metric used for checkpoint selection: the JuniperMapper Photo Interpretation validation split (67 images) and COCO mask AP$_{50}$, with the selected checkpoint and all other settings unchanged.

\subsection{LiDAR processing}

The UAV LiDAR point cloud was initially pre-processed using DJI Terra (v.5.3.5) software. Owing to data-storage constraints and to reduce the computational requirements of subsequent processing steps, a subset corresponding to approximately one-sixth of the original point cloud was retained for further analysis. The remaining processing was conducted in R Statistical Software (v4.6.1; \citealp{rsoftware}) using the lidR R package (v.4.3.2; \citealp{lidr2020}, \citealp{lidr2026}). All the scripts to conduct the following processing operations are inside the \textit{R} folder in the repository.

First, Statistical Outlier Removal (SOR) operation was applied using 10 nearest neighbors and a three-standard-deviation rejection threshold. Then, a ground classification was performed. Three ground-classification algorithms were evaluated: Multiscale Curvature Classification (MCC), Progressive Morphological Filtering (PMF), and Progressive TIN Densification (PTD). Ground-classification performance was assessed qualitatively using elevation profiles along four representative transects and quantitatively against surveyed GNSS elevations. The last dataset was collected by a Leica GNSS receiver in the center of each plot with a centimetric accuracy.

For quantitative assessment, LiDAR-derived digital terrain models (DTMs) for each classification method of each GNSS reference point were extracted. Each DTM was built with a TIN-interpolated method from each classified point cloud on the \qty{0.05}{\metre} orthomosaic grid. Each DTM was then sampled with bilinear interpolation at the GNSS targets. Elevation residuals were calculated relative to the surveyed orthometric elevation, and the mean RMSE across the GNSS targets was used to select the ground-classification method subsequently used for point-cloud normalization.

The most accurate ground-classified point cloud was used to normalize heights of all the points subtracting from the ground height, producing a Canopy Height Model (CHM). The resulting CHM had a spatial resolution of \qty{0.25}{\metre} and was aligned to the orthomosaic grid. This resolution was chosen because at \qty{5}{\centi\metre} only \qty{5.5}{\percent} of the cells contained a direct LiDAR return, compared with \qty{69}{\percent} at \qty{0.25}{\metre}. Cells without returns were retained as NoData rather than interpolated.

\subsection{LiDAR-based refinement}

The CHM was used as a post-processing to refine the shrub maps. It was not supplied as an input to any deep-learning model, so the same refinement could be applied independently to the predictions of both architectures. Because CHM heights close to the ground were within the LiDAR noise, the CHM was used only to separate trees from shrubs; predictions dominated by bare ground were not filtered and were left to the model confidence threshold.

The refinement process was carried out following the next steps. First, the CHM cells whose centers fell inside a predicted polygon were extracted. From them, the number of valid cells, the coverage (proportion of cells with a valid height), the mean, maximum and 90th-percentile height (P90), and the tall-cell fraction ($f$) were calculated. The tall-cell fraction was the proportion of valid cells exceeding a height threshold $h$.

Each prediction was then assigned to one of four post-processing categories, evaluated in the following order:

\begin{itemize}
\item Insufficient support. Predictions with a coverage below 0.3 or fewer than five valid cells were retained without modification. Predictions smaller than one CHM cell or outside the CHM extent have no valid cells and fall into this category.
\item Rejected. Predictions with P90 $> h$ and a tall-cell fraction $\geq f$ were removed as trees.  
\item Clipped. Predictions with at least one cell above $h$ that did not meet the rejection rule had their tall cells removed, and the largest remaining connected part was retained. If that part was smaller than \qty{0.25}{\square\metre}, the prediction was rejected.  
\item Unchanged. All other predictions were retained as they were.  
\end{itemize}

The thresholds $h$ and $f$ were calibrated on the 79 study area vegetation masks, corresponding with 58 shrubs and 21 trees. The growth form (tree or shrub) of each mask was assigned from the species identified in the field. Each mask representing \textit{Q. ilex}, \textit{P. pinaster} and \textit{C. monogyna} was defined as tree. The rest were shrubs.

All combinations of $h \in \{1.3, 1.5, 1.7, 1.9\}$~\unit{\metre} and $f \in \{0, 0.1, 0.2, 0.3, 0.5\}$ were applied to the masks with the procedure above. The selection rule was defined before calibration. Among the combinations that neither rejected nor clipped any shrub mask, the one that rejected most tree masks was chosen, with ties resolved in favor of the higher $h$ and then the higher $f$.

Because the thresholds were selected on the same masks, the selection was repeated in a leave-one-ROI-out procedure. Each mask was assigned to the \qty{9}{\square\metre} field ROI it overlapped most, or to the nearest ROI if it overlapped none (three masks, at most \qty{0.85}{\metre} away). For each of the 26 ROIs, the thresholds were selected on the masks of the remaining ROIs and applied to the held-out masks.

The selected parameters were then applied unchanged to the predictions of both architectures. Performance before and after LiDAR refinement was compared using the same independent study-area validation procedure.

\section{Results}

\subsection{Model development and selection}

The Mask R-CNN model was evaluated on the PI validation subset seven times during the 4,000 training iterations (Table~\ref{tab:maskrcnn_iterations}). Mask AP$_{50}$ increased from \qty{78.9}{\percent} to a maximum of \qty{84.5}{\percent}, and mask AP$_{50:95}$ from \qty{52.4}{\percent} to \qty{56.5}{\percent}. Mask AP$_{50}$ then decreased slightly and stabilized at around \qty{83}{\percent} until the end of training. The epoch 6 with the highest mask AP$_{50}$ was therefore retained. The highest mask AP$_{50:95}$ differed from that of the selected checkpoint by only 0.2 percentage points, so the secondary criterion was not needed.

\begin{table}[ht]
    \centering
    \caption{
        Mask AP$_{50}$ and mask AP$_{50:95}$ on the validation set during Mask R-CNN training. The model was evaluated every 2 epochs (570 iterations; 570 training images, batch size 2). The best value of each metric is shown in bold. The checkpoint with the highest mask AP$_{50}$ (epoch 6) was kept as the final model.
    }
    \label{tab:maskrcnn_iterations}
    \begin{tabular}{rrcc}
        \toprule
        Epoch & Iteration & Mask AP$_{50}$ (\%) & Mask AP$_{50:95}$ (\%) \\
        \midrule
        2 &  570 & 78.9          & 52.4 \\
        4 & 1140 & 83.0          & 55.5 \\
        6 & 1710 & \textbf{84.5} & 56.3 \\
        8 & 2280 & 84.2          & \textbf{56.5} \\
        10 & 2850 & 83.0          & 56.1 \\
        12 & 3420 & 83.2          & 56.0 \\
        14 & 3990 & 83.1          & 56.0 \\
        \bottomrule
    \end{tabular}
\end{table}

Domain adaptation of SAM 3 was also monitored on the PI validation subset using mask AP$_{50}$ and mask AP$_{50:95}$. Validation performance improved rapidly during the first epochs and subsequently reached a plateau (Table~\ref{tab:sam3_epochs}). Mask AP$_{50}$ increased from \qty{85.2}{\percent} at epoch 1 to \qty{88.5}{\percent} at epoch 4, while mask AP$_{50:95}$ increased from \qty{61.3}{\percent} to \qty{63.8}{\percent} over the same interval. From epochs 4 to 18, performance remained comparatively stable, with mask AP$_{50}$ ranging from \qty{87.7}{\percent} to \qty{88.4}{\percent} and mask AP$_{50:95}$ from \qty{63.5}{\percent} to \qty{64.7}{\percent}.

The highest validation performance was obtained at epoch 19, with a mask AP$_{50}$ of \qty{89.0}{\percent} and a mask AP$_{50:95}$ of \qty{64.8}{\percent}, both the highest values in their respective metrics during training. Consequently, epoch 19 was retained as the final domain-adapted SAM 3 checkpoint.

\begin{table}[ht]
    \centering
    \caption{
        Assessed SAM 3 epochs on the PI-validation set.
    }
    \label{tab:sam3_epochs}
    \begin{tabular}{ccc@{\qquad}ccc}
        \toprule
        Epoch & Mask AP$_{50}$ (\%) & Mask AP$_{50:95}$ (\%) & Epoch & Mask AP$_{50}$ (\%) & Mask AP$_{50:95}$ (\%) \\
        \midrule
        1 & 85.2 & 61.3 & 11 & 87.7 & 64.4 \\
        2 & 86.7 & 63.1 & 12 & 88.1 & 64.2 \\
        3 & 87.5 & 63.4 & 13 & 88.1 & 64.3 \\
        4 & 88.5 & 63.8 & 14 & 88.1 & 64.4 \\
        5 & 88.2 & 64.3 & 15 & 87.7 & 63.9 \\
        6 & 88.4 & 64.1 & 16 & 88.3 & 64.7 \\
        7 & 87.9 & 64.2 & 17 & 88.0 & 64.4 \\
        8 & 88.0 & 64.1 & 18 & 88.2 & 64.5 \\
        9 & 87.7 & 63.5 & \textbf{19} & \textbf{89.0} & \textbf{64.8} \\
        10 & 88.2 & 64.4 & 20--30 & -- & -- \\
        \bottomrule
    \end{tabular}
\end{table}

Although the training schedule was configured for 30 epochs, optimization became unstable during epoch 20, when the matching cost became non-finite. Subsequent batches were skipped by the trainer. Accordingly, the zero mask AP values recorded for epochs 20-30 do not represent a genuine deterioration in model performance, but rather the absence of effective parameter updates following numerical failure. Epoch 19 therefore represents both the highest-performing validation checkpoint and the last checkpoint obtained before training instability.

\subsection{Exploratory SAM 3 prompt analysis}

On the Photo Interpretation validation split, the text prompt changed SAM 3's accuracy only slightly unless it added a color or condition modifier (Table~\ref{tab:prompt_sensitivity}). The predefined prompt ``shrub'' reached a mask AP$_{50}$ of \qty{89.0}{\percent} and a mask AP$_{50:95}$ of \qty{64.8}{\percent}, the same values recorded when the checkpoint was selected. Six of the alternative prompts (``low shrub'', ``bush'', ``Mediterranean shrub'', ``small bush'', ``woody shrub'' and ``small shrub'') scored within 0.6 percentage points of mask AP$_{50}$ of ``shrub'' (\qtyrange{88.5}{88.9}{\percent}). Their mask AP$_{50:95}$ was equally close (\qtyrange{64.3}{65.1}{\percent}). No single prompt ranked first on both metrics: ``shrub'' had the highest mask AP$_{50}$ and ``woody shrub'' the highest mask AP$_{50:95}$, and both leads are smaller than 0.4 percentage points. On a 67-image validation set, such differences do not support preferring any synonym over the predefined prompt.

Prompts that added a color or condition modifier performed consistently worse. Compared with ``shrub'', mask AP$_{50}$ fell by 1.8 percentage points for ``dry shrub'', by 2.7 for ``burned shrub'' and by 10.2 for ``brown shrub''. ``Brown shrub'' also had the lowest mask AP$_{50:95}$ (\qty{56.1}{\percent}, against \qty{64.8}{\percent} for ``shrub'').

Because the alternative prompts were evaluated on the same validation split used to select the SAM 3 checkpoint with the ``shrub'' prompt, the analysis was treated as exploratory rather than as a post-hoc prompt optimization. The predefined prompt ``shrub'' was therefore retained for the principal inter-model comparison.

\begin{table}[ht]
    \centering
    \caption{
        SAM 3 prompt sensitivity on the JuniperMapper Photo Interpretation validation split (67 images), using the domain-adapted checkpoint selected with the prompt ``shrub''. COCO mask metrics; prompts are ordered by mask AP$_{50}$.
    }
    \label{tab:prompt_sensitivity}
    \begin{tabular}{lcc}
        \toprule
        Prompt & Mask AP$_{50}$ (\%) & Mask AP$_{50:95}$ (\%) \\
        \midrule
        shrub & 89.0 & 64.8 \\
        low shrub & 88.9 & 64.9 \\
        bush & 88.8 & 64.9 \\
        Mediterranean shrub & 88.8 & 65.0 \\
        small bush & 88.6 & 64.5 \\
        woody shrub & 88.6 & 65.1 \\
        small shrub & 88.5 & 64.3 \\
        dry shrub & 87.3 & 63.7 \\
        burned shrub & 86.4 & 63.9 \\
        brown shrub & 78.9 & 56.1 \\
        \bottomrule
    \end{tabular}
\end{table}

\subsection{Held-out JuniperMapper evaluation}

After model selection, the fixed Mask R-CNN and domain-adapted SAM 3 models were evaluated on the FW and PI test datasets. The FW results are presented first because the deployment operating points were selected on them.

\subsubsection{Field-work dataset performance}

Regarding F1 response to the prediction-confidence threshold on the FW dataset (Figure~\ref{fig:internal_performance_fw}), the threshold differed markedly between the two architectures.

At IoU~=~0.50 and S-IoU~=~0.50, Mask R-CNN rose steeply from F1 = \qty{15}{\percent} (IoU) and \qty{17}{\percent} (S-IoU) at a confidence of \qty{0}{\percent}, where all 31,744 candidate predictions were retained, to a broad plateau. F1 peaked at \qty{35}{\percent} confidence (IoU \qty{70.4}{\percent}; S-IoU \qty{75.8}{\percent}) and remained within 1 percentage point of this maximum between \qty{20}{\percent} and \qty{45}{\percent}. Above \qty{60}{\percent} confidence it declined gradually, and above \qty{85}{\percent} more steeply, reaching approximately \qty{46}{\percent} at \qty{95}{\percent}.

The domain-adapted SAM 3 model showed a unimodal, sharply peaked response. F1 increased progressively from \qty{22}{\percent} at \qty{0}{\percent} confidence, where precision was low because of the large number of low-scoring predictions, to a maximum at \qtyrange{55}{60}{\percent} (IoU \qty{74.9}{\percent}; S-IoU \qty{78.3}{\percent}). It then fell rapidly to \qty{60}{\percent} at \qty{80}{\percent} confidence and \qty{5}{\percent} at \qty{95}{\percent}, because few predictions received scores above \qty{85}{\percent}.

Sensitivity to the choice of operating point therefore differed between models. On the S-IoU~=~0.50 curve, Mask R-CNN was robust to a poorly chosen threshold. From its optimum prediction-confidence threshold (\qty{35}{\percent}) to \qty{55}{\percent}, 20 points off, F1 dropped by only 2 points, from \qty{75.8}{\percent} to \qty{73.8}{\percent}. SAM 3 was fragile, because moving 20 points off its optimum at \qty{60}{\percent} prediction confidence cost a drop of \qty{16.2}{\percent} in F1.

\begin{figure}[ht]
    \centering
    \includegraphics[width=1\textwidth]{figures/fig_models_internal_performance_fw.png}
    \caption{
        Mask R-CNN and SAM 3 performance evaluated on FW set using IoU and S-IoU metrics across different prediction-confidence score thresholds ($X$ axis).
    }
    \label{fig:internal_performance_fw}
\end{figure}

At the stricter IoU~=~0.75 and S-IoU~=~0.75, F1 decreased for both models, but the shape of each response was largely preserved. SAM 3 again achieved higher values (S-IoU \qty{68.4}{\percent} at \qty{60}{\percent} confidence; IoU \qty{53.6}{\percent} at \qty{65}{\percent}) than Mask R-CNN (S-IoU \qty{64.2}{\percent} at \qty{35}{\percent}; IoU \qty{45.3}{\percent} at \qty{80}{\percent}). For Mask R-CNN, IoU-based F1 increased slightly with confidence. This indicates that higher-scoring predictions delineated crown boundaries more accurately.

Comparing the two matching metrics, S-IoU-based F1 exceeded IoU-based F1 at every confidence threshold for both models. The difference was moderate at a threshold of 0.50 (approximately 3-5 percentage points) but increased to 15-20 points at 0.75.

The operating-point selection was done by applying the selection criterion, i.e., the highest confidence threshold whose F1 at S-IoU~=~0.50 was within 1 percentage point of the maximum. This gave operating points of \qty{60}{\percent} for SAM 3 and \qty{45}{\percent} for Mask R-CNN (Table~\ref{tab:fw_operating_points}).

\begin{table}[ht]
    \centering
    \caption{
        Precision, recall and F1 of each model at its selected operating point on the FW dataset (1,771 reference shrubs), matched with S-IoU~=~0.50.
    }
    \label{tab:fw_operating_points}
    \begin{tabular}{lcccc}
        \toprule
        Model & Threshold & Precision & Recall & F1 \\
        \midrule
        SAM 3 & \qty{60}{\percent} & \qty{81.6}{\percent} & \qty{75.3}{\percent} & \qty{78.3}{\percent} \\
        Mask R-CNN & \qty{45}{\percent} & \qty{82.9}{\percent} & \qty{68.8}{\percent} & \qty{75.2}{\percent} \\
        \bottomrule
    \end{tabular}
\end{table}

For SAM 3, the criterion selected the maximum itself (\qty{78.3}{\percent}), because F1 fell to \qty{77.1}{\percent} at \qty{65}{\percent}. For Mask R-CNN, the maximum at \qty{35}{\percent} (\qty{75.8}{\percent}) lay within a broad plateau, and the criterion selected \qty{45}{\percent} (\qty{75.2}{\percent}). This raised precision from \qty{79.6}{\percent} to \qty{82.9}{\percent} at the cost of a reduction in recall from \qty{72.4}{\percent} to \qty{68.8}{\percent}, and reduced the number of predictions from 1,553 to 1,413.

At these operating points, SAM 3 achieved the higher F1 (\qty{78.3}{\percent}; precision \qty{81.6}{\percent}, recall \qty{75.3}{\percent}), followed by Mask R-CNN (\qty{75.2}{\percent}; precision \qty{82.9}{\percent}, recall \qty{68.8}{\percent}). The ranking was consistent across all four evaluation criteria. Mask R-CNN achieved the higher precision, whereas the advantage of SAM 3 resulted primarily from its higher recall.

\subsubsection{PI test performance}

On the PI test dataset, F1 was consistently higher than on the FW dataset, and both models reached their maxima at higher confidence thresholds (Figure~\ref{fig:internal_performance_pi}). At IoU~=~0.50 and S-IoU~=~0.50, Mask R-CNN rose steeply from F1 = \qty{14}{\percent} (IoU) and \qty{15}{\percent} (S-IoU) at \qty{0}{\percent} confidence, where all 19,200 candidate predictions were retained, to \qty{65}{\percent} and \qty{67}{\percent} at \qty{5}{\percent}. F1 then increased steadily to a maximum at \qty{80}{\percent} (IoU \qty{84.8}{\percent}; S-IoU \qty{88.6}{\percent}), stayed within 1 percentage point of this maximum between \qty{65}{\percent} and \qty{85}{\percent}, and declined above \qty{90}{\percent}, reaching \qtyrange{76}{79}{\percent} at \qty{95}{\percent}. SAM 3 again showed an unimodal response. F1 increased from \qty{18}{\percent} at \qty{0}{\percent} confidence to a maximum at \qtyrange{65}{70}{\percent} (IoU \qty{88.5}{\percent}; S-IoU \qty{91.1}{\percent}), then declined sharply to \qty{86.0}{\percent} at \qty{80}{\percent}, \qty{55.8}{\percent} at \qty{90}{\percent} and \qty{14.1}{\percent} at \qty{95}{\percent}, because few predictions received scores above \qty{85}{\percent}. As on the FW dataset, SAM 3 was the more sensitive to the choice of operating point. A departure of 10 percentage points from the optimum reduced its F1 by up to 5.1 points, compared with up to 2.3 points for Mask R-CNN. A departure of 20 points reduced SAM 3's F1 by up to 35.3 points, compared with 1.2 points for Mask R-CNN.

\begin{figure}[ht]
    \centering
    \includegraphics[width=1\textwidth]{figures/fig_models_internal_performance_pi.png}
    \caption{
        Mask R-CNN and SAM 3 performance evaluated on PI test set using IoU and S-IoU metrics across different prediction-confidence score thresholds ($X$ axis).
    }
    \label{fig:internal_performance_pi}
\end{figure}

At IoU~=~0.75 and S-IoU~=~0.75, F1 on the PI dataset decreased much less than on the FW dataset, and the shape of each response was preserved. SAM 3 again achieved higher values (S-IoU \qty{87.6}{\percent} at \qty{70}{\percent} confidence; IoU \qty{80.5}{\percent} at \qty{75}{\percent}) than Mask R-CNN (S-IoU \qty{85.0}{\percent} at \qty{80}{\percent}; IoU \qty{72.3}{\percent} at \qty{85}{\percent}). For both models, the best F1 at IoU~=~0.75 occurred at a confidence equal to or higher than the best at IoU~=~0.50. This again indicates that higher-scoring predictions delineated crown boundaries more accurately.

S-IoU-based F1 exceeded IoU-based F1 at every confidence threshold for both models, but the differences were smaller than on the FW dataset. At the respective maxima, the difference was 2.6-3.8 percentage points at a threshold of 0.50 and 7.1-12.7 points at 0.75, compared with 15-20 points on the FW dataset.

Applying the same selection criterion to the PI dataset gave operating points of \qty{70}{\percent} for SAM 3 and \qty{85}{\percent} for Mask R-CNN (Table~\ref{tab:pi_operating_points}). For SAM 3, the criterion selected the maximum itself (\qty{91.1}{\percent}; precision \qty{90.4}{\percent}, recall \qty{91.9}{\percent}). For Mask R-CNN, it selected \qty{85}{\percent} (\qty{88.2}{\percent}; precision \qty{92.4}{\percent}, recall \qty{84.4}{\percent}) instead of the \qty{80}{\percent} maximum (\qty{88.6}{\percent}). The ranking of the two models, where SAM 3 was higher than Mask R-CNN, was the same as on the FW dataset and was consistent across all four evaluation criteria.

\begin{table}[ht]
    \centering
    \caption{
        Precision, recall and F1 of each model at the operating point selected on the PI test dataset (690 reference shrubs), matched with S-IoU~=~0.50.
    }
    \label{tab:pi_operating_points}
    \begin{tabular}{lcccc}
        \toprule
        Model & Threshold & Precision & Recall & F1 \\
        \midrule
        SAM 3 & \qty{70}{\percent} & \qty{90.4}{\percent} & \qty{91.9}{\percent} & \qty{91.1}{\percent} \\
        Mask R-CNN & \qty{85}{\percent} & \qty{92.4}{\percent} & \qty{84.4}{\percent} & \qty{88.2}{\percent} \\
        \bottomrule
    \end{tabular}
\end{table}

The thresholds selected on the PI dataset were 10 (SAM 3) and 40 (Mask R-CNN) percentage points higher than those selected on the FW dataset. Both models therefore assigned lower confidence scores to shrubs in the field-based reference data. When the FW thresholds were applied to the PI dataset, F1 at S-IoU~=~0.50 fell by 2.2 points for SAM 3 (\qty{88.9}{\percent}) and 2.8 points for Mask R-CNN (\qty{85.4}{\percent}), mainly through lower precision, while recall remained at \qtyrange{90}{96}{\percent}. When the PI thresholds were applied to the FW dataset, F1 fell by 3.8 points for SAM 3 (\qty{74.5}{\percent}) and 13.1 points for Mask R-CNN (\qty{62.1}{\percent}), mainly through lower recall. Operating points derived from the field-based data therefore transferred to the photo-interpretation domain with a much smaller loss of accuracy than the reverse. This supports the use of the FW-derived thresholds (\qty{60}{\percent} for SAM 3 and \qty{45}{\percent} for Mask R-CNN) for deployment.

\subsection{LiDAR ground classification}

Of the 37 targets, 26 were covered by all DTMs and were used for the comparison, so every method was evaluated on the same set of targets. PMF had the lowest vertical RMSE (\qty{0.612}{\metre}), followed by PTD with the \qty{30}{\degree}/\qty{2}{\metre} (\qty{0.641}{\metre}) and \qty{30}{\degree}/\qty{3}{\metre} (\qty{0.645}{\metre}) parameterizations, MCC (\qty{0.727}{\metre}) and PTD with \qty{40}{\degree}/\qty{3}{\metre} (\qty{0.787}{\metre}) (Table~\ref{tab:ground_classification}). All DTMs overestimated the GNSS heights, with mean errors between \qty{0.356}{\metre} (PMF) and \qty{0.517}{\metre} (PTD \qty{40}{\degree}/\qty{3}{\metre}). PMF also had the lowest share of targets with an absolute error above \qty{0.5}{\metre} (\qty{15.4}{\percent}), compared with \qty{23.1}{\percent} for both PTD \qty{30}{\degree} variants, \qty{30.8}{\percent} for PTD \qty{40}{\degree}/\qty{3}{\metre} and \qty{34.6}{\percent} for MCC. PMF was therefore retained for point-cloud normalization and CHM generation.

\begin{table}[ht]
    \centering
    \caption{
        Vertical accuracy of the DTMs derived from each ground-classification algorithm, assessed at 26 GNSS targets. Errors are computed as DTM minus GNSS orthometric height.
    }
    \label{tab:ground_classification}
    \resizebox{\textwidth}{!}{%
        \begin{tabular}{lccccc}
            \toprule
            Method & Mean error (\unit{\metre}) & SD (\unit{\metre}) & MAE (\unit{\metre}) & RMSE (\unit{\metre}) & Targets with $|\text{error}| > \qty{0.5}{\metre}$ (\unit{\percent}) \\
            \midrule
            PMF & 0.356 & 0.508 & 0.392 & 0.612 & 15.4 \\
            PTD \qty{30}{\degree}/\qty{2}{\metre} & 0.388 & 0.520 & 0.425 & 0.641 & 23.1 \\
            PTD \qty{30}{\degree}/\qty{3}{\metre} & 0.390 & 0.525 & 0.426 & 0.645 & 23.1 \\
            MCC & 0.500 & 0.539 & 0.524 & 0.727 & 34.6 \\
            PTD \qty{40}{\degree}/\qty{3}{\metre} & 0.517 & 0.605 & 0.554 & 0.787 & 30.8 \\
            \bottomrule
        \end{tabular}%
    }
\end{table}

The terrain profiles in Figure~\ref{fig:transect_profiles} show the classification outputs along four transects. MCC and PMF produced broadly similar ground surfaces. MCC kept ground returns beneath and between shrub canopies (e.g. Transect 4, \qtyrange{-10}{-7.5}{\metre}), whereas PMF left some ground returns in the non-ground class (e.g. Transect 4, around \qty{10}{\metre}). PTD produced the noisiest classification, with frequent non-ground labels on clearly bare terrain (e.g. Transects 2 and 3, \qtyrange{-5}{2}{\metre}).

\begin{figure}[ht]
    \centering
    \includegraphics[width=1\textwidth]{figures/transect_plot.png}
    \caption{
        Terrain profiles (distance along transect vs. elevation, in meters) for the MCC, PMF and PTD classifications. Each point is labeled as ground or other return. PTD is shown with its best parameterization: a maximum iteration angle of \qty{30}{\degree} and a maximum iteration distance of \qty{2}{\metre}, with a \qty{10}{\metre} seed resolution. Rows correspond to transects and columns to classification algorithms.
    }
    \label{fig:transect_profiles}
\end{figure}

\subsection{Calibration of the LiDAR refinement rule}

CHM measurements were available for all 79 field-reference masks. The median number of valid CHM cells per mask was 15 for shrubs and 96 for trees. With the minimum-support criteria, 10 shrub masks, all of them small \textit{L. stoechas}, had insufficient LiDAR support and would be retained without modification. No tree mask fell into this category.

Mean CHM height differed clearly between growth forms (Figure~\ref{fig:species_mean_height}). \textit{Q. ilex} (\qty{2.87}{\metre}) and \textit{P. pinaster} (\qty{2.45}{\metre}) were the tallest species. They were followed by \textit{C. monogyna} (\qty{1.00}{\metre}), whose mean height was closer to that of the tallest shrub, \textit{C. laurifolius} (\qty{0.67}{\metre}), than to the other tree species. \textit{E. vagans} (\qty{0.21}{\metre}) and \textit{L. stoechas} (\qty{0.13}{\metre}) were the lowest.

The species means hid an overlap at mask level. Six tree masks had mean heights below \qty{1}{\metre}: three \textit{Q. ilex}, two \textit{P. pinaster} and one \textit{C. monogyna}. That range was also reached by \textit{C. laurifolius} masks (up to \qty{1.03}{\metre}). The 90th-percentile height (P90) ranged from 0 to \qty{1.46}{\metre} in shrub masks and from 1.35 to \qty{6.75}{\metre} in tree masks. No CHM cell inside any shrub mask exceeded \qty{1.63}{\metre}. The CHM also contained negative values down to \qty{-0.05}{\metre}, a noise level close to the mean height of \textit{L. stoechas} (\qty{0.13}{\metre}).

\begin{figure}[ht]
    \centering
    \includegraphics[width=1\textwidth]{figures/fig_species_mean_height_points.png}
    \caption{
        Cell-weighted mean CHM height of the species with field-reference masks, colored by growth form.
    }
    \label{fig:species_mean_height}
\end{figure}

Table~\ref{tab:refinement_calibration} shows how the threshold combinations classified the field-reference masks. A tall-cell fraction of 0 produced the same results as 0.1 at every height, because P90 $> h$ already implies that about \qty{10}{\percent} of the cells exceed $h$. Fractions above 0.1 reduced the number of rejected trees without protecting any additional shrub.

Lowering $h$ increased tree rejection but started to affect shrubs. At \qty{1.5}{\metre}, 18 of 21 trees were rejected and no shrub was rejected, but six shrub masks were clipped, and the threshold was only \qty{4}{\centi\metre} above the highest shrub P90 (\textit{C. laurifolius}, \qty{1.46}{\metre}). At \qty{1.3}{\metre}, one shrub mask was rejected.

Applying the selection rule gave P90 $>$ \qty{1.7}{\metre} with a tall-cell fraction $\geq$ 0.1. This rejected 14 trees, clipped 4 and retained 3, without rejecting or clipping any shrub, since \qty{1.7}{\metre} lies above every shrub cell in the sample.

\begin{table}[ht]
    \centering
    \caption{
        Outcome of the refinement rule on the field-reference masks (21 trees, 58 shrubs) for each combination of P90 threshold and tall-cell fraction (minimum coverage 0.3). Ten shrub masks had insufficient LiDAR support in all combinations. The selected combination is in bold.
    }
    \label{tab:refinement_calibration}
    \resizebox{\textwidth}{!}{%
        \begin{tabular}{ccccccc}
            \toprule
            P90 $>$ (\unit{\metre}) & Tall fraction $\geq$ & Trees rejected & Trees clipped & Trees retained & Shrubs rejected & Shrubs clipped \\
            \midrule
            1.3 & 0.1 & 21 & 0 & 0 & 1 & 6 \\
            1.3 & 0.3 & 19 & 2 & 0 & 1 & 6 \\
            1.3 & 0.5 & 14 & 7 & 0 & 0 & 7 \\
            1.5 & 0.1 & 18 & 3 & 0 & 0 & 6 \\
            1.5 & 0.3 & 14 & 7 & 0 & 0 & 6 \\
            1.5 & 0.5 & 14 & 7 & 0 & 0 & 6 \\
            \textbf{1.7} & \textbf{0.1} & \textbf{14} & \textbf{4} & \textbf{3} & \textbf{0} & \textbf{0} \\
            1.7 & 0.3 & 13 & 5 & 3 & 0 & 0 \\
            1.7 & 0.5 & 12 & 6 & 3 & 0 & 0 \\
            1.9 & 0.1 & 13 & 2 & 6 & 0 & 0 \\
            1.9 & 0.3 & 12 & 3 & 6 & 0 & 0 \\
            1.9 & 0.5 & 11 & 4 & 6 & 0 & 0 \\
            \bottomrule
        \end{tabular}%
    }
\end{table}

With a minimum coverage of 0.5, two \textit{P. pinaster} masks (coverage 0.44 and 0.48) would also have had insufficient support, and the selected combination would have rejected 13 trees instead of 14.

The three retained trees were one \textit{Q. ilex} (P90 = \qty{1.35}{\metre}), one \textit{P. pinaster} (\qty{1.44}{\metre}) and one \textit{C. monogyna} (\qty{1.37}{\metre}). None contained a cell above \qty{1.7}{\metre}, and their P90 values were below that of the tallest \textit{C. laurifolius}, so they could not be separated from shrubs by height alone. \textit{C. monogyna} also commonly grows as a tall shrub. The four clipped trees had P90 between 1.55 and \qty{1.69}{\metre}.

P90 was preferred over maximum height because it was robust to isolated returns. Two \textit{L. stoechas} masks had maximum heights of 1.52 and \qty{1.54}{\metre} but P90 of only 0.22 and \qty{0.33}{\metre}.

In the leave-one-ROI-out procedure, the selection rule chose \qty{1.7}{\metre} with a tall-cell fraction of 0.1 in 25 of the 26 folds. In the remaining fold, leaving out one tree made 1.7 and \qty{1.9}{\metre} equivalent on the training masks, and \qty{1.9}{\metre} was selected. Pooled over folds, the held-out masks gave 13 of 21 trees rejected, 4 clipped and 4 retained, with no shrub rejected or clipped. This is one tree fewer than in the calibration sample.

When clipping shrubs was allowed during selection, \qty{1.5}{\metre} was chosen in 25 folds. When the ROI containing the tallest \textit{C. laurifolius} was held out, the threshold dropped to \qty{1.3}{\metre} and that shrub was rejected. This indicates that a threshold placed above all shrub heights, rather than just above the highest shrub P90, generalizes better to unseen masks.

\subsection{Shrub maps of the study area}

Each model was deployed over the Encinacorba orthomosaic at its FW operating point (\qty{45}{\percent} for Mask R-CNN and \qty{60}{\percent} for SAM 3). After the tile predictions were merged, the Mask R-CNN map contained 1,726 shrub polygons covering \qty{7909}{\square\metre}, and the SAM 3 map contained 3,377 polygons covering \qty{15700}{\square\metre}. The median polygon area was \qty{2.5}{\square\metre} for Mask R-CNN and \qty{3.1}{\square\metre} for SAM 3. SAM 3 therefore predicted about twice as many objects and twice the area as Mask R-CNN.

\subsection{Cross-domain validation in the study area}

Both models transferred poorly to the field reference of the study area (Table~\ref{tab:field_validation}). Mask R-CNN matched 1 of the 58 reference shrubs (recall \qty{1.7}{\percent}) and SAM 3 matched 2 (recall \qty{3.4}{\percent}). F1 was \qty{3.2}{\percent} for Mask R-CNN and \qty{6.1}{\percent} for SAM 3. At the same operating points, the two models had reached \qty{75.2}{\percent} and \qty{78.3}{\percent} on the JuniperMapper FW dataset, although under the stricter S-IoU~=~0.50 criterion.

\begin{table}[ht]
    \centering
    \caption{
        Accuracy of the deployed shrub maps against the field reference (58 shrubs, 21 trees) before LiDAR refinement. Scored predictions are the true plus false positives. Mean IoU is calculated over the matched shrubs and is reported on a 0--1 scale. Trees matched is the number of reference trees matched by a prediction, and tree commission is that number as a percentage of the 21 reference trees.
    }
    \label{tab:field_validation}
    \resizebox{\textwidth}{!}{%
        \begin{tabular}{lcccccccccc}
            \toprule
            Model & Scored predictions & TP & FP & FN & Precision (\%) & Recall (\%) & F1 (\%) & Mean IoU & Trees matched & Tree commission (\%) \\
            \midrule
            Mask R-CNN & 4 & 1 & 3 & 57 & 25.0 & 1.7 & 3.2 & 0.106 & 6 & 28.6 \\
            SAM 3 & 8 & 2 & 6 & 56 & 25.0 & 3.4 & 6.1 & 0.852 & 13 & 61.9 \\
            \bottomrule
        \end{tabular}
    }
\end{table}

The matched shrubs differed in quality between models (Figure~\ref{fig:map_predictions}). The two SAM 3 matches were \textit{C. laurifolius} crowns of 4.0 and \qty{13.7}{\square\metre}, delineated with IoU values of 0.81 and 0.89 and confidence scores of \qty{88}{\percent} and \qty{95}{\percent}. Both predictions lay mostly outside the plots and were counted as true positives because matched predictions were scored wherever they lay. The single Mask R-CNN match was a \qty{5.1}{\square\metre} prediction centered on a \textit{C. monogyna} tree (IoU 0.49 with the tree polygon, first row of the Figure~\ref{fig:map_predictions}). It overlapped a \qty{0.8}{\square\metre} \textit{L. stoechas} polygon with an IoU of 0.106, just above the matching threshold of 0.10. It therefore met the matching criterion but did not delineate that shrub. Under a matching threshold of 0.25, Mask R-CNN would have had no true positive.

\begin{figure}[ht]
    \centering
    \includegraphics[width=1\textwidth]{figures/fig_map_plots_evaluation_nomatterport.png}
    \caption{
        Predicted shrubs by model compared with the reference data divided between shrub or tree mask. A representative validation plot is shown in each row, with its species composition and the occupation percentage of each one in parenthesis.
    }
    \label{fig:map_predictions}
\end{figure}

Detection was strongly species dependent. Neither model detected any of the 4 \textit{E. vagans} masks (Figure~\ref{fig:map_predictions}), and the 46 \textit{L. stoechas} masks, which made up \qty{79}{\percent} of the reference shrubs, were detected only through the marginal Mask R-CNN match described above. SAM 3 detected 2 of the 8 \textit{C. laurifolius} masks, and Mask R-CNN detected none.

All in-plot false positives of both models fell on trees. The three Mask R-CNN false positives overlapped \textit{Q. ilex} crowns. Of the six SAM 3 false positives, four overlapped \textit{Q. ilex}, one \textit{P. pinaster} and one \textit{C. monogyna}. Outside the plots, the predictions also matched many reference trees. Mask R-CNN matched 6 of the 21 trees (3 \textit{Q. ilex}, 2 \textit{P. pinaster} and 1 \textit{C. monogyna}; median IoU 0.39), and SAM 3 matched 13 (8 \textit{Q. ilex}, 4 \textit{P. pinaster} and 1 \textit{C. monogyna}; median IoU 0.55). SAM 3 delineated several tree crowns closely, with IoU values above 0.88 for four of them.

Precision was \qty{25}{\percent} for both models, but it was calculated from only 4 and 8 scored predictions. It is therefore too unstable to rank the models.

\subsection{Effect of LiDAR refinement}

Over the whole map, the refinement removed most of the predicted area (Table~\ref{tab:refinement_outcome}). For Mask R-CNN, 715 of the 1,726 predictions (\qty{41.4}{\percent}) were rejected as trees. They covered \qty{5319}{\square\metre}, or \qty{67.2}{\percent} of the mapped area. For SAM 3, 1,942 of the 3,377 predictions (\qty{57.5}{\percent}) were rejected, covering \qty{12053}{\square\metre} (\qty{76.8}{\percent}). The median P90 of the rejected predictions was \qty{2.48}{\metre} for Mask R-CNN and \qty{2.50}{\metre} for SAM 3, compared with \qty{0.79}{\metre} and \qty{1.04}{\metre} for the predictions left unchanged. A further 138 Mask R-CNN and 237 SAM 3 predictions were clipped. Clipping removed only \qty{32}{\square\metre} and \qty{44}{\square\metre}, respectively, and no clipped prediction fell below the \qty{0.25}{\square\metre} minimum. Of the remaining predictions, 130 (Mask R-CNN) and 205 (SAM 3) had insufficient LiDAR support and were retained without modification. After refinement, the Mask R-CNN map contained 1,011 polygons covering \qty{2559}{\square\metre}, and the SAM 3 map contained 1,435 polygons covering \qty{3603}{\square\metre}.

\begin{table}[ht]
    \centering
    \caption{
        Outcome of the LiDAR refinement on the deployed shrub maps.
    }
    \label{tab:refinement_outcome}
    \resizebox{\textwidth}{!}{%
        \begin{tabular}{lcccccc}
            \toprule
            Model & Predictions & Rejected (trees) & Clipped & Unchanged & Insufficient support & Retained area (\unit{\square\metre}) \\
            \midrule
            Mask R-CNN & 1,726 & 715 & 138 & 743 & 130 & 2,559 (of 7,909) \\
            SAM 3 & 3,377 & 1,942 & 237 & 993 & 205 & 3,603 (of 15,700) \\
            \bottomrule
        \end{tabular}%
    }
\end{table}

Against the field reference, the refinement removed tree commission without losing any matched shrub (Table~\ref{tab:refinement_validation}). The trees matched by predictions fell from 6 to 1 for Mask R-CNN and from 13 to 3 for SAM 3. The tree-commission rate fell from \qty{28.6}{\percent} to \qty{4.8}{\percent} and from \qty{61.9}{\percent} to \qty{14.3}{\percent}, respectively. All three in-plot Mask R-CNN false positives and four of the six SAM 3 false positives were rejected. Precision therefore rose from \qty{25}{\percent} to \qty{100}{\percent} for Mask R-CNN and from \qty{25}{\percent} to \qty{50}{\percent} for SAM 3. Recall did not change, so F1 increased only marginally, to \qty{3.4}{\percent} and \qty{6.5}{\percent}.

\begin{table}[ht]
    \centering
    \caption{
        Accuracy of the shrub maps against the field reference before and after LiDAR refinement.
    }
    \label{tab:refinement_validation}
    \resizebox{\textwidth}{!}{%
        \begin{tabular}{llccccccccc}
            \toprule
            Model & Stage & Scored predictions & TP & FP & FN & Precision (\%) & Recall (\%) & F1 (\%) & Trees matched & Tree commission (\%) \\
            \midrule
            Mask R-CNN & Before & 4 & 1 & 3 & 57 & 25.0 & 1.7 & 3.2 & 6 & 28.6 \\
            Mask R-CNN & After & 1 & 1 & 0 & 57 & 100.0 & 1.7 & 3.4 & 1 & 4.8 \\
            SAM 3 & Before & 8 & 2 & 6 & 56 & 25.0 & 3.4 & 6.1 & 13 & 61.9 \\
            SAM 3 & After & 4 & 2 & 2 & 56 & 50.0 & 3.4 & 6.5 & 3 & 14.3 \\
            \bottomrule
        \end{tabular}%
    }
\end{table}

The two SAM 3 \textit{C. laurifolius} matches were retained unchanged, with the same IoU values (Figure~\ref{fig:map_refined_predictions}). The single Mask R-CNN true positive, the prediction on the \textit{C. monogyna} tree, was clipped (P90 \qty{1.51}{\metre}) and kept. It still overlapped the \textit{L. stoechas} polygon with IoU 0.107. The Mask R-CNN precision of \qty{100}{\percent} therefore rests on one prediction that mainly covers a tree.

The trees that were still matched after refinement were those that the calibration had already shown to be inseparable from shrubs by height. Both models still matched \textit{C. monogyna} tree 44, whose reference mask had a P90 of \qty{1.37}{\metre} and was retained in the calibration. SAM 3 also still matched \textit{Q. ilex} tree 65 (reference P90 \qty{1.35}{\metre}, also retained in the calibration) with IoU 0.83, and \textit{Q. ilex} tree 58 through a \qty{4.5}{\square\metre} prediction with a P90 of \qty{0.96}{\metre}. The two SAM 3 false positives that remained in the plots were the predictions on \textit{C. monogyna} (P90 \qty{1.51}{\metre}) and on \textit{Q. ilex} tree 58. Neither contained any CHM cell above \qty{1.7}{\metre}, so both were left unchanged.

\begin{figure}[ht]
    \centering
    \includegraphics[width=1\textwidth]{figures/fig_map_plots_lidar_refinement_nomatterport.png}
    \caption{
        Refined shrubs maps by model compared with the reference data divided between shrub or tree mask. A representative validation plot is shown in each row, with its species composition and the occupation percentage of each one in parenthesis.
    }
    \label{fig:map_refined_predictions}
\end{figure}

\subsection{Qualitative error analysis}

Visual inspection identified two recurrent error types (Figure~\ref{fig:visual_inspection}). First, shadowed regions were occasionally classified as shrubs and were not removed by the LiDAR height filter because their estimated heights did not resemble tree crowns. Second, some Mask R-CNN masks showed straight boundaries corresponding to inference-tile edges. Comparable tile-boundary artifacts were not visually evident in the SAM 3 output.

\begin{figure}[ht]
    \centering
    \includegraphics[width=1\textwidth]{figures/noyet.png}
    \caption{
        Captured errors by a visual inspection, included shadowed regions (column A) and straight boundaries (column B) for the two models.
    }
    \label{fig:visual_inspection}
\end{figure}

\section{Discussion}

The current trained Mask R-CNN model differed slightly to the one trained in \citet{khaldi2024individual}. Although both models share hyperparameters and training strategy, 
Comparing the best F1 achieved on each dataset (Table A1), our model was within 1.7 percentage points of the original on the PI test set across all four evaluation criteria (e.g. IoU~=~0.50: \qty{84.8}{\percent} vs \qty{84.84}{\percent}; S-IoU~=~0.50: \qty{88.6}{\percent} vs \qty{87.87}{\percent}). On the FW dataset it was 1.1–2.6 points lower at a threshold of 0.50 and at IoU~=~0.75, and 4.1 points lower for S-IoU~=~0.75 (\qty{64.2}{\percent} vs \qty{68.32}{\percent}). At the original study's confidence thresholds (\qty{90}{\percent} for PI and \qty{50}{\percent} for FW), our model was 1.5–4.4 points lower. It retained fewer predictions (543 vs 646 on PI and 1,345 vs 1,706 on FW), with higher precision and lower recall. Our model therefore assigned lower confidence scores than the original, which shifted its FW optimum from \qty{50}{\percent} to \qty{35}{\percent}, while its achievable segmentation performance remained comparable. This difference in score behavior between two training runs of the same architecture supports deriving operating points separately for each model rather than adopting a fixed value from previous work. The confidence-threshold responses were also qualitatively consistent with the original study. On the PI data, F1 increased with confidence and peaked at a high threshold (\qty{80}{\percent} vs \qty{90}{\percent}). On the FW data, F1 rose rapidly and then formed a broad plateau, although ours extended over a narrower range (\qtyrange{20}{45}{\percent} vs \qtyrange{20}{85}{\percent}). Both studies showed a similar decrease from PI to FW data at a threshold of 0.50 (14.4 and 12.8 points for IoU and S-IoU, vs 12.45 and 11.01). This indicates that the domain shift between photo-interpreted and field-based references is reproducible. At a threshold of 0.75, however, the S-IoU decrease was larger in our study (20.8 vs 16.5 points), suggesting that our model fragmented or merged field-observed crowns more often. The domain-adapted SAM 3 model exceeded the original JuniperMapper performance under seven of the eight criteria, by 1.4–6.6 points, and matched it for FW S-IoU~=~0.75 (\qty{68.4}{\percent} vs \qty{68.32}{\percent}). The largest gains were for IoU~=~0.75 (PI \qty{80.5}{\percent} vs \qty{73.95}{\percent}; FW \qty{53.6}{\percent} vs \qty{47.91}{\percent}), indicating more accurate crown delineation rather than only more detections.

However, overall the IoU and S-IoU metrics of both models follow a similar shape, with S-IoU's F1-score consistently higher than IoU's.

the Mask R-CNN model trained in this study closely matched the performance reported for the original JuniperMapper model (Khaldi et al., 2024).

S-IoU credits a prediction against the union of all overlapping ground-truth instances rather than requiring a single one-to-one match, so fragmented or merged canopy predictions still count as (partial) successes. This directional definition also means that, for S-IoU, the true positives entering the precision calculation are not the same instances as those entering the recall calculation: precision uses prediction-directed matches (each prediction scored against its unioned overlapping ground truths), while recall and its false negatives are computed independently, label-directed (each ground truth scored against its unioned overlapping predictions). As a result, TP + FN does not equal the total number of ground-truth instances for the S-IoU rows below, unlike the IoU rows, where one-to-one matching guarantees the two are consistent. The IoU-based F1 slightly increase for Mask R-CNN with confidence could indicate that higher-scoring predictions delineated crown boundaries more accurately. The S-IoU-based F1 exceeded IoU-based F1 at every confidence threshold for both models, and this could suggest that a substantial share of errors under the stricter criterion arose from fragmented or merged crowns and boundary misplacement rather than from undetected shrubs. The IoU and S-IoU curves converged at high confidence thresholds. Fragmentation and merging errors were therefore concentrated among lower-confidence predictions, whereas the most confident predictions generally corresponded to single, well-delineated crowns.

SAM 3 Zero-Shot BAD FOR DOMAIN SPECIFIC TASKS. Zero-shot mask AP falls from 56.4 (COCO) to 15.2 on a 100-domain specialized benchmark (Dabaja and Çelík 2026), i.e. zero-shot text-prompt segmentation degrades sharply outside SAM 3's pretraining distribution and cannot be assumed to transfer to vegetation/remote-sensing imagery without domain-specific evaluation. In an orchard fruit-annotation study (Wang, 2026), SAM 3 reached higher recall than SAM (\qty{91}{\percent} vs. \qty{81}{\percent}) but markedly lower precision (\qty{29}{\percent} vs. \qty{42}{\percent}) — i.e. SAM 3's dominant failure mode is over-segmentation / false positives, not missed detections. In remote sensing, text prompts inject a ground-level semantic bias that degrades coordinate regression on overhead imagery, while visual (exemplar) prompts align well to overhead geometry (Dabaja and Çelík 2026). Our imagery (Zenmuse L1 orthomosaics, \qty{5}{\centi\metre} GSD) is exactly this overhead case.

S-IoU-based F1 exceeded IoU-based F1 at every confidence threshold for both models, but the differences in PI were smaller than on the FW dataset. This could indicate that fragmentation, merging and boundary errors had less influence on performance within the photo-interpretation domain.

SAM 3 has evaluated internally with shrub prompt, but then several one has been tested. There could be a slightly decrease in performance because the confidence value to select masks has not been set programmatically with each prompt. Maybe account that in further research.

Notably, the best domain-adapted SAM 3 checkpoints modestly outperform the Mask R-CNN reproduction (Section 4.1) at matched IoU~=~0.50 and S-IoU~=~0.50 operating points on both test sets: \qty{88.5}{\percent} vs. \qty{84.8}{\percent} F1 (IoU~=~0.50) on PI, and \qty{74.9}{\percent} vs. \qty{70.4}{\percent} on FW. This suggests that, despite starting from a much weaker zero-shot baseline, lightweight frozen-backbone decoder fine-tuning lets SAM 3 match or exceed a fully fine-tuned Mask R-CNN on this dataset — provided the checkpoint is taken before the numerical instability identified above (epoch 19 was the last checkpoint before it).



The models mostly predict taller vegetation and miss the low \textit{Lavandula}. That suggested a recall and domain problem, and CHM filtering cannot fix it. The three masks in the validation set corresponding to a tree species but with low heigh values may correspond to low resprouts after the 2003 fire.

Higher overlap of the SAM 3 model could be one of the reason of the no presence of border effect between tiles.

All ground-classification algorithms placed the terrain surface above the GNSS reference: 25 of the 26 targets had a positive error for every method. This consistent positive bias suggests that low vegetation, such as herbaceous cover, small shrubs and litter, was partly kept in the ground class. Dense low vegetation can stop the laser before it reaches the soil, so the lowest available return marks the top of that layer rather than the ground surface. The more conservative filters reduced this effect. PMF had the smallest bias (\qty{0.356}{\metre}) and MCC and PTD \qty{40}{\degree}/\qty{3}{\metre} the largest (\qty{0.500}{\metre} and \qty{0.517}{\metre}), which matches their different tolerance for low objects close to the terrain. The largest single error (about \qty{2.6}{\metre} for every method) occurred at target 31, located at the center of a \textit{Q. ilex} crown. There the dense evergreen canopy blocked the laser from reaching the soil, so none of the algorithms could find reliable ground points beneath it. The interpolated DTM was then controlled by returns from the crown or its edges. Excluding this target lowers the RMSE to \qty{0.346}{\metre} for PMF, \qty{0.416}{\metre} for PTD \qty{30}{\degree}/\qty{2}{\metre}, \qty{0.530}{\metre} for MCC and \qty{0.618}{\metre} for PTD \qty{40}{\degree}/\qty{3}{\metre}, and the mean error to \qtyrange{0.267}{0.436}{\metre}. The ranking of the methods stays the same. Part of the remaining bias may also come from a systematic vertical offset between the GNSS and LiDAR height references, which this design cannot separate from the vegetation effect.

Picos et al. (2020), for example, used \qty{0.25}{\metre} for the ground–shrub division and 2 and \qty{3}{\metre} for shrub–stem separation in two UAV-LiDAR \textit{Eucalyptus} plantation plots. However, this work only serves as a justification of LiDAR refinement operation. The thresholds come from a different system, and is not a safe start here. The median valid cells per mask is 15 for shrubs and 96 for trees. That is still small for tiny \textit{L. stoechas} masks, but usable. Three trees are only \texttt{keep}. Ids 65 (\textit{Q. ilex}, \texttt{h\_p90} of 1.35), 60 (\textit{P. pinaster}, 1.44) and 44 (\textit{C. monogyna}, 1.37) have a \texttt{tall\_fraction} of 0. Their heights overlap \textit{C. laurifolius} (up to 1.46), so height alone cannot remove them. Trees 44 and 65 were also the trees still matched by predictions after the refinement of the deployed maps. These could be low resprouts after the 2003 fire. The models mostly predict taller vegetation and miss the low \textit{Lavandula}. That is a recall and domain problem, and CHM filtering cannot fix it.

Precision rests on 1 to 8 scored predictions, so the values aren't reliable enough to rank the models. The Mask R-CNN precision of \qty{100}{\percent} after refinement rests on a single prediction that mainly covers a \textit{C. monogyna} tree and matches a \textit{L. stoechas} polygon with an IoU of only 0.107. Because shrub masks are not well predicted across all plots, the LiDAR refinement worked to remove the tree commission but not to improve the masks of the few well detected shrubs. The LiDAR is not proposed for the proposal of new shrub locations because the thresholds are validated with the annotations, so there is no way to assess the gains of this process without a bias.

Each model is deployed with its own checkpoint and its own confidence threshold, \qty{45}{\percent} for Mask R-CNN and \qty{60}{\percent} for SAM 3, and not with a single common value. The scores of the two architectures are not calibrated on the same scale, so a shared threshold would not make the comparison fairer. The consequence is that the deployment results reflect the operating point of each model as well as the architecture itself: a lower threshold keeps more low-confidence masks, which increases the number of predictions and the chance of a match at the cost of precision, and a higher one does the opposite.

\section{Conclusions}

\begin{enumerate}
\item The refinement is a tree-commission fix. It removes most masks that land on trees (tree commission down \qtyrange{77}{83}{\percent}) without losing any matched shrub, so it is safe to apply. Precision rises only because false positives are removed from a very small set of predictions, and F1 stays below \qty{10}{\percent} because recall does not change.  
\item More CHM filtering won't help. Tightening the height rule would only start removing true shrubs (\textit{Cistus} and \textit{Crataegus} reach \qty{1.5}{\metre}), and the FPs mostly sit in the \qtyrange{0.3}{1.7}{\metre} band.  
\item The bottleneck is the models. Recall of \qtyrange{1.7}{3.4}{\percent} means they hit 1 to 2 of the 58 annotated shrubs. \textbf{Check whether these are matching problems (IoU 0.10 is already lenient) or species/height mismatches}.  
\item \textbf{Look and describe the remaining FPs}. Overlay the retained SAM 3 masks on the orthomosaic to see what they are. If many are real shrubs that were never digitized, the ground truth is the limit, not the model. Mask R-CNN's only hit is a matching artifact, and SAM 3's two hits are real \textit{C. laurifolius} crowns.
\end{enumerate}

\bibliographystyle{plainnat}
\bibliography{references}


\section{Appendix I}

\begin{table}[ht]
    \centering
    \caption{
        Comparison with Khaldi et al. (2024), Table 4 (F1, \unit{\percent}).
    }
    \label{tab:khaldi_comparison}
    \begin{tabular}{llccc}
        \toprule
        Dataset & Metric & Khaldi et al. & Mask R-CNN (this study) & SAM 3 \\
        \midrule
        PI & IoU~=~0.50 & 84.84 & 84.8 / 82.7 & 88.5 \\
        PI & S-IoU~=~0.50 & 87.87 & 88.6 / 86.3 & 91.1 \\
        PI & IoU~=~0.75 & 73.95 & 72.3 / 72.0 & 80.5 \\
        PI & S-IoU~=~0.75 & 84.77 & 85.0 / 83.3 & 87.6 \\
        \midrule
        FW & IoU~=~0.50 & 72.39 & 70.4 / 69.4 & 74.9 \\
        FW & S-IoU~=~0.50 & 76.86 & 75.8 / 74.3 & 78.3 \\
        FW & IoU~=~0.75 & 47.91 & 45.3 / 44.7 & 53.6 \\
        FW & S-IoU~=~0.75 & 68.32 & 64.2 / 63.9 & 68.4 \\
        \bottomrule
    \end{tabular}
\end{table}

\end{document}
